\documentclass[11pt]{article}
\usepackage[margin=1in]{geometry}
\usepackage{amsmath,amssymb,amsthm}
\usepackage{booktabs}
\usepackage{graphicx}
\usepackage[hidelinks]{hyperref}
\usepackage{microtype}

\newtheorem{theorem}{Theorem}
\newcommand{\D}{\mathbb D}
\newcommand{\T}{\mathbb T}
\newcommand{\C}{\mathbb C}

\title{\textbf{A Counterexample to the Hardy-Space Norm Conjecture}\\
\large Conjecture 4.3 of arXiv:2511.08000}
\author{Automated mathematical discovery run \texttt{fa\_banach\_001}\\
Agent \texttt{agent\_lane\_08}}
\date{23 July 2026}

\begin{document}
\maketitle

\begin{center}
\fbox{\parbox{0.91\textwidth}{
\textbf{Status: candidate counterexample, likely valid.}
The argument is exact: it uses finite rational coefficient sums and no
numerical integration.
}}
\end{center}

\section{Source and target}

The source is C.~B\'en\'eteau, R.~Cheng, C.~Felder, D.~Khavinson,
M.~Manolaki, and K.~Maronikolakis, \emph{Metric projections, zeros of
optimal polynomial approximants, and some extremal problems in Hardy
spaces}, arXiv:2511.08000v2, Conjecture~4.3 on page~12
\cite{source}.

\begin{figure}[ht]
\centering
\includegraphics[width=\textwidth]{figures/open_problem_crop.png}
\caption{Conjecture 4.3 in the source paper, page 12.}
\end{figure}

With normalized Lebesgue measure \(m\) on \(\T\), the conjecture reads:

\begin{quote}
Let \(1<p<\infty\) and \(f\in H^p\) with \(f(0)\ne0\).  If \(J\) is any
nonconstant inner function, then there exists \(c=c_f\in\T\) such that
\[
        \lVert 1-Jf\rVert_p>\lVert 1-cf\rVert_p.
\]
\end{quote}

\section{Idea of the proof}

The exponent \(p=4\) makes the proposed inequality completely algebraic:
for every analytic \(h\in H^4\), Parseval gives
\(\lVert h\rVert_4^4=\lVert h^2\rVert_2^2\).  We choose a short real
polynomial \(f\).  The quantity \(\lVert1-cf\rVert_4^4\), for
\(|c|=1\), then becomes a quadratic polynomial in \(x=\operatorname{Re}c\);
its global minimum is read off on \([-1,1]\).  For the left side we choose a
single Blaschke factor \(J\).  Multiplication by \(J\) is isometric in
\(H^2\), so only one finite coefficient pairing remains.  The two exact
rational values occur in the strict reverse order.

\section{Counterexample}

\begin{theorem}
Set
\[
 p=4,\qquad
 f(z)=\frac12-3z+2z^3-z^5,\qquad
 J(z)=\frac{z-\frac14}{1-\frac14z}.
\]
Then \(f(0)=\frac12\), and \(J\) is a nonconstant inner function with
\(J(0)=-\frac14\).  Moreover,
\[
 \lVert1-Jf\rVert_4^4=\frac{1344927}{4096}
 <\frac{5697}{16}
 =\min_{|c|=1}\lVert1-cf\rVert_4^4.
\]
Consequently, every \(c\in\T\) satisfies
\[
       \lVert1-Jf\rVert_4<\lVert1-cf\rVert_4,
\]
so Conjecture~4.3 is false.
\end{theorem}

\begin{proof}
All inner products and norms below are taken on \(\T\) with normalized
Lebesgue measure.  If \(h(z)=\sum_{n\ge0}h_nz^n\in H^4\), then
\begin{equation}\label{eq:parseval}
 \lVert h\rVert_4^4
 =\int_{\T}|h|^4\,dm
 =\lVert h^2\rVert_2^2
 =\sum_{n\ge0}\bigl|\widehat{h^2}(n)\bigr|^2.
\end{equation}

\paragraph{The constant-phase side.}
For \(|c|=1\), direct expansion gives the nonzero coefficients of
\((1-cf)^2\):
\[
\begin{array}{c|rrrrrrrrr}
n&0&1&2&3&4&5&6&8&10\\ \hline
[z^n](1-cf)^2&
\frac{(c-2)^2}{4}&-3c(c-2)&9c^2&2c(c-2)&-12c^2&
-c(c-2)&10c^2&-4c^2&c^2 .
\end{array}
\]
Applying \eqref{eq:parseval} and writing \(c=e^{i\theta}\) and
\(x=\cos\theta\), we obtain
\begin{align}
\lVert1-cf\rVert_4^4
 &=\frac{6625}{16}
   +\frac12\cos(2\theta)-\frac{117}{2}\cos\theta \notag\\
 &=\frac{6617}{16}+x^2-\frac{117}{2}x. \label{eq:right}
\end{align}
The derivative \(2x-\frac{117}{2}\) is negative throughout
\([-1,1]\).  Thus \eqref{eq:right} is minimized at \(x=1\), and
\begin{equation}\label{eq:rightmin}
 \min_{|c|=1}\lVert1-cf\rVert_4^4
 =\frac{6617}{16}+1-\frac{117}{2}
 =\frac{5697}{16}.
\end{equation}

\paragraph{The inner-function side.}
Put \(u=Jf\).  Expanding \((1-u)^2=1-2u+u^2\) in \(H^2\) gives
\begin{equation}\label{eq:leftidentity}
\lVert1-Jf\rVert_4^4
=1+4\lVert f\rVert_2^2+\lVert f^2\rVert_2^2
-4\operatorname{Re}\bigl((Jf)(0)\bigr)
+2\operatorname{Re}\bigl((Jf)(0)^2\bigr)
-4\operatorname{Re}\langle f,Jf^2\rangle .
\end{equation}
Here we used that multiplication by the inner function \(J\) is an
isometry on \(H^2\), and hence
\[
\langle Jf,J^2f^2\rangle=\langle f,Jf^2\rangle.
\]
Finite coefficient sums yield
\[
 \lVert f\rVert_2^2=\frac{57}{4},
 \qquad
 \lVert f^2\rVert_2^2=\frac{5697}{16}.
\]
For completeness, the Taylor polynomial through degree five of \(Jf^2\) is
\[
 Jf^2=-\frac1{16}+\frac{63}{64}z-\frac{1281}{256}z^2
 +\frac{7423}{1024}z^3+\frac{27903}{4096}z^4
 -\frac{164609}{16384}z^5+O(z^6).
\]
Since \(f\) has degree five, this is enough to compute the mixed term:
\[
 \langle f,Jf^2\rangle
 =\frac{353249}{16384}.
\]
Finally, \((Jf)(0)=-\frac18\).  Substitution into
\eqref{eq:leftidentity} gives
\begin{align*}
\lVert1-Jf\rVert_4^4
&=\frac{6625}{16}+\frac12+\frac1{32}
  -4\left(\frac{353249}{16384}\right)\\
&=\frac{1344927}{4096}.
\end{align*}
Comparison with \eqref{eq:rightmin} is strict:
\[
 \frac{5697}{16}-\frac{1344927}{4096}
 =\frac{113505}{4096}>0.
\]
Taking fourth roots preserves the inequality and proves the claim.
\end{proof}

\section{Verification}

The script \texttt{code/verify\_counterexample.py} recomputes every displayed
quantity using Python's exact \texttt{Fraction} arithmetic.  It uses the
Taylor formula
\[
[z^0]J=-a,\qquad [z^n]J=(1-a^2)a^{n-1}\quad(n\ge1),
\qquad a=\frac14,
\]
and polynomial convolution; it performs no quadrature or floating-point
comparison.  The command
\[
\texttt{python code/verify\_counterexample.py}
\]
returns the two exact fourth powers, their positive gap, and a passing status.

\section{Scope and novelty check}

This example disproves Conjecture~4.3 exactly, and it also satisfies the
stronger side condition \(J(0)\ne0\) appearing in the preceding
Proposition~4.2 of the source.  It does \emph{not} by itself construct an
optimal polynomial approximant with a zero in \(\D\): Proposition~4.2 is a
one-way sufficient mechanism, so the broader zero-free OPA question is not
settled here.

Before promotion, the run's lightweight indexes were searched for
arXiv:2511.08000, the exact title, the Hardy-space norm inequality, inner
functions, and OPA zeros.  Bounded web searches on 23 July 2026 checked the
current arXiv v2 record, the Canadian Journal of Mathematics publication
record, the exact conjecture phrase and label, and close
Hardy-space/optimal-polynomial-approximant terms.  These searches found the
source and background OPA papers but no later proof, counterexample,
correction, or explicit resolution of Conjecture~4.3.  Novelty is therefore
plausible but remains subject to expert literature review.

\section{Human-review recommendation}

The counterexample is recommended for mathematical review as \emph{likely
valid}.  The most useful manual check is the mixed coefficient pairing
\(\langle f,Jf^2\rangle=353249/16384\); the exact verifier independently
reconstructs it from polynomial and Blaschke coefficients.

\subsection*{Exact verifier report}

Running the bundled verifier in the packet directory returns:
\begin{verbatim}
||1-Jf||_4^4 = 1344927/4096
inf_|c|=1 ||1-cf||_4^4 = 5697/16
strict reversed gap = 113505/4096
PASS: Conjecture 4.3 is violated for every unimodular c.
\end{verbatim}

\subsection*{Suggested review checks}

\begin{enumerate}
\item Confirm that \(J(z)=(z-\frac14)/(1-\frac14z)\) is inner and
      \(J(0)\ne0\).
\item Re-sum the nine coefficients of \((1-cf)^2\) and the quadratic
      in \(x=\operatorname{Re}c\).
\item Recompute the six Taylor coefficients of \(Jf^2\) used in the
      finite inner product.
\item Note that the conclusion is the strict reverse for every \(c\in\T\),
      so no issue of phase selection remains.
\end{enumerate}

\begin{thebibliography}{1}
\bibitem{source}
C.~B\'en\'eteau, R.~Cheng, C.~Felder, D.~Khavinson, M.~Manolaki, and
K.~Maronikolakis,
\emph{Metric projections, zeros of optimal polynomial approximants, and
some extremal problems in Hardy spaces},
arXiv:2511.08000v2, 2026.
\end{thebibliography}

\end{document}
